\scanpage{49}
\mspage{12}
\[
(14)\qquad\alpha:\Delta_{r'}\longrightarrow\Delta_r
\]
is given by
\[
\alpha^*(x,i_*)=(x,\underbrace{\alpha^*(i_*)}_{\substack{\Vert\,\mathrm{def}\\i_*\circ\alpha}})
\]
i.e. it is induced by the product semi-
simplicial structure in the second member of
\[
(15)\qquad X_{\Delta_*}\subset X\times I_{\Delta_*}.
\]
\uline{Remarks}\unsure{.} a) All these constructions generalize to the
case where one would replace \inserted{the ordered set} $I$ by a category \inserted{$\unsure{\mathcal{C}|C}$}, and
the map $X\mapsto X_i$ by a functor $\unsure{\mathcal{C}|C}\longrightarrow\text{category}$
of closed subsets of $X$ and of their inclusions
...\unsure{.|-}

b) If one starts from a family $(X_i)_{i\in I}$, \inserted{$I$ a set} and nothing
more, one can introduce ad hoc an order
relation on $I$ by \deleted{$i\leq j\unsure{=|\Longleftrightarrow}$}
\[
i\leq j\ \overset{\mathrm{def}}{\Longleftrightarrow}\ X_i\leq X_j.
\]
But in the presentation adopted, it is not
necessary to suppose that \inserted{the \unsure{increasing|increasing} map} $I\longrightarrow\text{set of subsets}$
closed in $X$ is injective, and still less
that it is an embedding of ordered sets,
i.e. that one has
\[
X_i\leq X_j\ \Longrightarrow\ i\leq j
\]
which is a converse of (1' b)).
\margin{48}
